% GENERATED FILE. Edit canonical lessons, then run npm run generate:books.
% canonical-source-hash: fnv1a64:4e1e4f2c63a7499b
% canonical-lessons: MW-C130-borlo, MW-C130-nath, MW-C130-payal, MW-C130-bajro, MW-C130-rabri

\chapter{Forehead ornament, Nose ring, Anklet, Millet, Millet porridge}
\label{ch:mw-forehead-ornament-nose-ring-anklet-millet-millet-porridge}
\hlchaptermodality{130}

\begin{chapteropening}
\textbf{By the end of this chapter:} \emph{I can say a forehead ornament, a nose ring, an anklet, millet and a millet porridge.}

Complete the last lesson of chapter 130: I can say a forehead ornament, a nose ring, an anklet, millet and a millet porridge.
\end{chapteropening}

\section[\texorpdfstring{\mw{बोरलो} (borlo)}{बोरलो (borlo)}]{\mw{बोरलो} (borlo) — a forehead ornament}

\label{lesson:MW-C130-borlo}



\begin{quote}
Before the new one: say the Marwadi for a thick bangle, then the Marwadi for a bangle.
\end{quote}

\subsection*{You'll want to know: \mw{बोरलो}}
\textbf{\mw{बोरलो}} — \emph{borlo} — \textquotedblleft{}a forehead ornament\textquotedblright{}.

\begin{cousinweb}
The origin of \textbf{\mw{बोरलो}} (\emph{borlo}) is not settled. CDIAL, searched by form and by meaning, has no entry for it.
\end{cousinweb}

\begin{grammarlens}[title={What you've built}]
One more word for this chapter.
\end{grammarlens}

\subsection*{Your turn}
\begin{itemize}
\raggedright
  \item \emph{Say it:} \emph{borlo}
  \item \emph{Say it:} \emph{borlo}, once more
  \item \emph{Say it:} say \emph{borlo}
  \item \emph{Recall:} say the Marwadi for a thick bangle, then the Marwadi for a bangle, then say \emph{borlo} again
\end{itemize}

\begin{tcolorbox}[breakable,colback=teal!4,colframe=teal!35!black,title={Before you move on}]
What does \mw{बोरलो} mean? (\textquotedblleft{}A forehead ornament\textquotedblright{}.) Say it once more.
\end{tcolorbox}
\section[\texorpdfstring{\mw{नथ} (nath)}{नथ (nath)}]{\mw{नथ} (nath) — a nose ring}

\label{lesson:MW-C130-nath}



\begin{quote}
Before the new one: say the Marwadi for a bangle, then the Marwadi for a forehead ornament.
\end{quote}

\subsection*{You'll want to know: \mw{नथ}}
\textbf{\mw{नथ}} — \emph{nath} — \textquotedblleft{}a nose ring\textquotedblright{}.

\begin{cousinweb}
\textbf{\mw{नथ}} (\emph{nath}) is a tadbhava, a word inherited from Sanskrit and reshaped by sound change. It comes from \emph{nasta}, \textquotedblleft{}nose\textquotedblright{}. CDIAL gives no Marwari form in this entry (CDIAL 7031); the match is through Gujarati \emph{nath}. CDIAL lists Gujarati \emph{nath}, \textquotedblleft{}nose ring for women\textquotedblright{}, and Hindi \emph{nāth}, \textquotedblleft{}nose rope\textquotedblright{}, under it.
\end{cousinweb}

\begin{grammarlens}[title={What you've built}]
One more word for this chapter.
\end{grammarlens}

\subsection*{Your turn}
\begin{itemize}
\raggedright
  \item \emph{Say it:} \emph{nath}
  \item \emph{Say it:} \emph{nath}, once more
  \item \emph{Say it:} say \emph{nath}
  \item \emph{Recall:} say the Marwadi for a bangle, then the Marwadi for a forehead ornament, then say \emph{nath} again
\end{itemize}

\begin{tcolorbox}[breakable,colback=teal!4,colframe=teal!35!black,title={Before you move on}]
What does \mw{नथ} mean? (\textquotedblleft{}A nose ring\textquotedblright{}.) Say it once more.
\end{tcolorbox}
\section[\texorpdfstring{\mw{पायल} (pāyal)}{पायल (pāyal)}]{\mw{पायल} (pāyal) — an anklet}

\label{lesson:MW-C130-payal}



\begin{quote}
Before the new one: say the Marwadi for a forehead ornament, then the Marwadi for a nose ring.
\end{quote}

\subsection*{You'll want to know: \mw{पायल}}
\textbf{\mw{पायल}} — \emph{pāyal} — \textquotedblleft{}an anklet\textquotedblright{}.

\begin{cousinweb}
\textbf{\mw{पायल}} (\emph{pāyal}) is a tadbhava, a word inherited and reshaped by sound change. It comes from \emph{*pādala}, \textquotedblleft{}appertaining to the feet\textquotedblright{}, a form CDIAL reconstructs from its descendants. CDIAL gives no Marwari form in this entry (CDIAL 8066); the match is through Hindi \emph{pāyal}, though CDIAL glosses that form only \textquotedblleft{}sure-footed\textquotedblright{}, not \textquotedblleft{}anklet\textquotedblright{}. It is built on \emph{pāda}, \textquotedblleft{}foot\textquotedblright{}.
\end{cousinweb}

\begin{grammarlens}[title={What you've built}]
One more word for this chapter.
\end{grammarlens}

\subsection*{Your turn}
\begin{itemize}
\raggedright
  \item \emph{Say it:} \emph{pāyal}
  \item \emph{Say it:} \emph{pāyal}, once more
  \item \emph{Say it:} say \emph{pāyal}
  \item \emph{Recall:} say the Marwadi for a forehead ornament, then the Marwadi for a nose ring, then say \emph{pāyal} again
\end{itemize}

\begin{tcolorbox}[breakable,colback=teal!4,colframe=teal!35!black,title={Before you move on}]
What does \mw{पायल} mean? (\textquotedblleft{}An anklet\textquotedblright{}.) Say it once more.
\end{tcolorbox}
\section[\texorpdfstring{\mw{बाजरो} (bājro)}{बाजरो (bājro)}]{\mw{बाजरो} (bājro) — millet}

\label{lesson:MW-C130-bajro}



\begin{quote}
Before the new one: say the Marwadi for a nose ring, then the Marwadi for an anklet.
\end{quote}

\subsection*{You'll want to know: \mw{बाजरो}}
\textbf{\mw{बाजरो}} — \emph{bājro} — \textquotedblleft{}millet\textquotedblright{}.

\begin{cousinweb}
\textbf{\mw{बाजरो}} (\emph{bājro}) is a tadbhava, a word inherited and reshaped by sound change. It comes from \emph{*bājjara}, \textquotedblleft{}millet\textquotedblright{}, a form CDIAL reconstructs from its descendants. CDIAL gives no Marwari form in this entry (CDIAL 9201); the match is through Hindi \emph{bājrā} and Gujarati \emph{bājro}.
\end{cousinweb}

\begin{grammarlens}[title={What you've built}]
One more word for this chapter.
\end{grammarlens}

\subsection*{Your turn}
\begin{itemize}
\raggedright
  \item \emph{Say it:} \emph{bājro}
  \item \emph{Say it:} \emph{bājro}, once more
  \item \emph{Say it:} say \emph{bājro}
  \item \emph{Recall:} say the Marwadi for a nose ring, then the Marwadi for an anklet, then say \emph{bājro} again
\end{itemize}

\begin{tcolorbox}[breakable,colback=teal!4,colframe=teal!35!black,title={Before you move on}]
What does \mw{बाजरो} mean? (\textquotedblleft{}Millet\textquotedblright{}.) Say it once more.
\end{tcolorbox}
\section[\texorpdfstring{\mw{राबड़ी} (rābṛī)}{राबड़ी (rābṛī)}]{\mw{राबड़ी} (rābṛī) — a millet porridge}

\label{lesson:MW-C130-rabri}



\begin{quote}
Before the new one: say the Marwadi for an anklet, then the Marwadi for millet.
\end{quote}

\subsection*{You'll want to know: \mw{राबड़ी}}
\textbf{\mw{राबड़ी}} — \emph{rābṛī} — \textquotedblleft{}a millet porridge\textquotedblright{}.

\begin{cousinweb}
\textbf{\mw{राबड़ी}} (\emph{rābṛī}) is a tadbhava, a word inherited and reshaped by sound change. It comes from \emph{*rabbā}, \textquotedblleft{}gruel, thickened juice\textquotedblright{}, a form CDIAL reconstructs from its descendants. CDIAL gives no Marwari form in this entry (CDIAL 10623); the match is through Hindi \emph{rābṛī} and Gujarati \emph{rābṛī}. The \emph{ṛ} is an added suffix.
\end{cousinweb}

\begin{grammarlens}[title={What you've built}]
One more word for this chapter.
\end{grammarlens}

\subsection*{Your turn}
\begin{itemize}
\raggedright
  \item \emph{Say it:} \emph{rābṛī}
  \item \emph{Say it:} \emph{rābṛī}, once more
  \item \emph{Say it:} say \emph{rābṛī}
  \item \emph{Recall:} say the Marwadi for an anklet, then the Marwadi for millet, then say \emph{rābṛī} again
\end{itemize}

\begin{tcolorbox}[breakable,colback=teal!4,colframe=teal!35!black,title={Before you move on}]
What does \mw{राबड़ी} mean? (\textquotedblleft{}A millet porridge\textquotedblright{}.) Say it once more.
\end{tcolorbox}
